\makeatletter
\def\input@path{{./}{../}{../../}{../../../}{../../../../}{preamble/}{../preamble/}{../../preamble/}{../../../preamble/}}
\makeatother
\input{preamble_exams.tex}
\input{graphs/highlights.tex}
% \GTenHighlightsfalse % Uncomment to disable instructional highlights.
\input{graphs/specific_highlights.tex}

\hypersetup{pdftitle={Grade 10 Global Economics: Term 3, Lesson 1, C5},pdfsubject={Home practice: Standard Normal Areas and Critical Values}}
\title{Grade 10 Global Economics\\Term 3, Lesson 1: Home Practice\\C5 --- Standard Normal Areas and Critical Values}
\author{}
\date{}
\begin{document}
\maketitle
\setcounter{tocdepth}{1}
\setcounter{secnumdepth}{2}
\phantomsection\label{toc}
\tableofcontents
\section*{Evaluated criteria}
\begin{description}
\item[\textbf{C5}] Applies the standard normal distribution \(N(0,1)\) and standard normal tables to calculate probabilities, critical values, and probability areas in finance and economics contexts.
\end{description}
\textbf{How to practise.} Sketch the requested area before reading the table. Use the supplied left-cumulative entries, symmetry and complements; no external table is required. Try each task on paper before opening its solution link. Record the model, working, units and an interpretation; use the arrows to return.
All datasets and economic or financial models in this activity are illustrative, rather than observed market evidence.

\clearpage
\section{Introduction}\label{sec:introduction}
\subsection*{Highlight legend}
\phantomsection\label{tab:highlight-legend}
\begin{tabularx}{\linewidth}{@{}>{\raggedright\arraybackslash}p{.20\linewidth}>{\raggedright\arraybackslash}p{.22\linewidth}>{\raggedright\arraybackslash}X@{}}
\toprule
\textbf{Symbol or notation} & \textbf{Concept} & \textbf{Brief meaning}\\
\midrule
$\GTenHLValue{Z}\sim N(0,1)$ & \GTenHLValue{Standard normal} & Normal model on a standard scale. \\
$\HLCumulativeProbability{\Phi(z)}$ & \HLCumulativeProbability{Cumulative area} & Area to the left of z. \\
$\HLTailArea{1-\Phi(z)}$ & \HLTailArea{Upper tail} & Area to the right of z. \\
$\HLCriticalValue{z_p}$ & \HLCriticalValue{Critical value} & Position with left area p. \\
\bottomrule
\end{tabularx}
\subsection{Introductory example 1: A standardised demand indicator}\label{ex:intro-one}
An illustrative demand indicator $\GTenHLValue{Z}\sim N(0,1)$ is already on a standard scale with \HLPopulationMean{mean} $\HLPopulationMean{0}$ and \HLPopulationSD{standard deviation} $\HLPopulationSD{1}$. A positive score is above average. The \HLCumulativeProbability{left-cumulative area} $\HLCumulativeProbability{\Phi(1.20)=0.8849}$ means about $88.49\%$ of scores lie below 1.20.
\subsection{Introductory example 2: An unusually high finance indicator}\label{ex:intro-two}
A standardised financing-cost indicator is flagged above $\HLCriticalValue{c}$. If the flag rate is $\HLTailArea{0.05}$, the cumulative target is $\HLCumulativeProbability{0.95}$. The cutoff $c$ is a \HLCriticalValue{critical value}, a location on the horizontal axis, whereas $0.95$ is an area.
\section{Notation Reference and Practice Procedure}\label{sec:notation-and-procedure}
\subsection{Notation reference}\label{tab:notation-reference}
\begin{tabularx}{\linewidth}{@{}>{\raggedright\arraybackslash}p{.16\linewidth}>{\raggedright\arraybackslash}p{.18\linewidth}>{\raggedright\arraybackslash}p{.27\linewidth}>{\raggedright\arraybackslash}X@{}}
\toprule
\textbf{Symbol or notation} & \textbf{Concept} & \textbf{Definition} & \textbf{Meaning in the introductory examples}\\
\midrule
$\GTenHLValue{Z}\sim N(0,1)$ & \GTenHLValue{Standard normal} & Bell-shaped symmetric distribution with mean zero and variance one; $Z$ is unitless. & \textbullet\ \hyperref[ex:intro-one]{Example 1}: Demand indicator is already standardised. \\
$\HLCumulativeProbability{\Phi(z)}$ & \HLCumulativeProbability{Cumulative area} & $Phi(z)=P(Zle z)$; table convention used here. & \textbullet\ \hyperref[ex:intro-one]{Example 1}: 0.8849 is the area left of 1.20. \\
$\HLTailArea{1-\Phi(z)}$ & \HLTailArea{Upper tail} & Complement of the left-cumulative area. & \textbullet\ \hyperref[ex:intro-two]{Example 2}: A high-cost flag uses an upper tail. \\
$\HLCriticalValue{z_p}$ & \HLCriticalValue{Critical value} & The value satisfying $Phi(z_p)=p$. & \textbullet\ \hyperref[ex:intro-two]{Example 2}: The flag boundary has left area 0.95. \\
\bottomrule
\end{tabularx}
\subsection{Practice procedure}
\begin{enumerate}
\item Identify whether the task gives a cutoff and asks for an area or gives an area and asks for a cutoff.
\item Sketch a standard bell curve centred at zero. Mark bounds and decide whether the desired region is left, right, between or outside.
\item Read $\Phi(z)$ from the excerpt. Use $\Phi(-z)=1-\Phi(z)$, right area $1-\Phi(z)$ and band area $\Phi(b)-\Phi(a)$.
\item For inverse lookup, convert the stated tail or central area to a left-cumulative target, find its matching cutoff and interpret location separately from area.
\end{enumerate}
\subsection{Self-contained standard-normal table excerpt}
Every entry below is the area \textbf{to the left}: $\Phi(z)=P(Z\le z)$. Negative entries follow $\Phi(-z)=1-\Phi(z)$. In a full table, $z=1.20$ is read at row $1.2$, column $.00$; $z=0.85$ is row $0.8$, column $.05$. The three-decimal critical values below are supplied interpolated reference values.
\[
\begin{array}{c|rrrrrrrrr}
z&0&0.35&0.85&0.95&1.15&1.20&1.40&1.645&1.96\\\hline
\Phi(z)&.5000&.6368&.8023&.8289&.8749&.8849&.9192&.9500&.9750
\end{array}
\]
Table-based results are approximate; equality at one point has probability zero.
\section{General Worked Problems}\label{sec:general-worked}
\SectionProblemsTableOfContents{
\SectionProblemLink{prob:sec:general-worked0}{Lower and upper demand tails}\\
\SectionProblemLink{prob:sec:general-worked1}{Comparing two reporting bands}\\
\SectionProblemLink{prob:sec:general-worked2}{Choosing a standard-scale critical value}}

\subsection{Lower and upper demand tails}\label{prob:sec:general-worked0}
A retailer's standardised demand indicator $Z$ is modeled as $N(0,1)$. A low-demand alert is triggered below $-1.20$, and a separate reporting rule records scores below $0.85$.
\QuestionsTasks[sec:general-worked]{{5}/{sol:sec:general-worked0}/{Find each alert/reporting probability and the probability of a score above 0.85.}}
\subsubsection*{Worked solution}\label{sol:sec:general-worked0}
Use the standard normal because the indicator is already standardised. The table convention gives
\[
P(Z<0.85)=\HLCumulativeProbability{0.8023}.
\]
By symmetry $\Phi(-1.20)=1-\Phi(1.20)=1-0.8849=0.1151$. Thus the low-demand alert probability is $\HLTailArea{0.1151}$, or $11.51\%$. The upper tail above $0.85$ is
\[
P(Z>0.85)=1-0.8023=\HLTailArea{0.1977},
\]
or $19.77\%$. The score $0.85$ is not itself a probability.
\secInicio[prob:sec:general-worked0][sec:general-worked]


\subsection{Comparing two reporting bands}\label{prob:sec:general-worked1}
For a standardised economic indicator $Z\sim N(0,1)$, report A accepts $0.35<Z<1.40$; report B accepts $-0.95<Z<1.15$.
\QuestionsTasks[sec:general-worked]{{5}/{sol:sec:general-worked1}/{Find and compare the accepted areas. Find the probability outside report B's band.}}
\subsubsection*{Worked solution}\label{sol:sec:general-worked1}
Each band is a difference of left areas. For A,
\[
P(0.35<Z<1.40)=0.9192-0.6368=\HLProbability{0.2824}.
\]
For B first reflect the negative bound:
\[
\Phi(-0.95)=1-0.8289=0.1711,
\]
\[
P(-0.95<Z<1.15)=0.8749-0.1711=\HLProbability{0.7038}.
\]
B accepts $70.38\%$ compared with A's $28.24\%$, a difference of $42.14$ percentage points. Its outside area is $1-0.7038=\HLTailArea{0.2962}$. A band farther from zero can have less probability even when its width looks substantial: the centre of the normal curve is taller.
\secInicio[prob:sec:general-worked1][sec:general-worked]


\subsection{Choosing a standard-scale critical value}\label{prob:sec:general-worked2}
A bank flags unusually high standardised borrowing-cost scores $Z\sim N(0,1)$. Policy A flags the upper $5\%$. Policy B flags values outside a symmetric central $95\%$ band.
\QuestionsTasks[sec:general-worked]{{5}/{sol:sec:general-worked2}/{Find the cutoff for A and the two cutoffs for B, explaining how the same percentage leads to different critical values.}}
\subsubsection*{Worked solution}\label{sol:sec:general-worked2}
For A, $P(Z>c)=0.05$ gives $\Phi(c)=0.95$. The excerpt supplies $\HLCriticalValue{c=1.645}$. For B the remaining $0.05$ is split into two tails of $0.025$, so
\[
\Phi(c_{\rm upper})=0.975,\qquad \Phi(c_{\rm lower})=0.025.
\]
The upper cutoff is $\HLCriticalValue{1.96}$ and symmetry gives $\HLCriticalValue{-1.96}$ for the lower one. A has one $5\%$ upper tail; B has two $2.5\%$ tails. Its more extreme positive cutoff leaves less area in that single upper tail.
\secInicio[prob:sec:general-worked2][sec:general-worked]
\section{Advanced Problems}\label{sec:advanced}
\SectionProblemsTableOfContents{
\SectionProblemLink{prob:sec:advanced0}{Two alert designs with equal total probability}}

\subsection{Two alert designs with equal total probability}\label{prob:sec:advanced0}
For $Z\sim N(0,1)$, policy A flags $Z>1.645$. Policy B flags $Z<-1.96$ or $Z>1.96$. A manager says both policies identify the same economic situations because both flag $5\%$.
\QuestionsTasks[sec:advanced]{{5}/{sol:sec:advanced0}/{Find the total flagged probability for each policy and the probability of being flagged by both. Evaluate the manager's claim.}}
\subsubsection*{Worked solution}\label{sol:sec:advanced0}
From the table, A flags $1-0.9500=0.0500$. B flags
\[
\Phi(-1.96)+1-\Phi(1.96)=0.0250+0.0250=0.0500.
\]
To be flagged by both, a score must exceed $1.96$: B's negative tail is outside A. Therefore
\[
P(\text{both})=1-\Phi(1.96)=\HLProbability{0.0250}.
\]
Only $2.5\%$ of all scores lie in the overlap. A also flags the moderate positive band $1.645<Z\le1.96$, while B also flags the negative tail below $-1.96$. Equal total probability does not mean equal events. The choice depends on whether only high-cost situations or unusually low and high situations matter.
\secInicio[prob:sec:advanced0][sec:advanced]

\end{document}

